O relatório da aula de Análise e Modelagem de Sistemas aponta baixo engajamento e atenção fragmentada. A sessão teve 120 minutos, com taxa de foco geral de 56,6\% e tempo médio de foco por aluno de apenas 13,7 minutos, sugerindo dificuldade em manter a atenção. A presença foi relativamente alta (82,1\%), mas o engajamento foi desigual, com apenas um aluno mantendo foco contínuo de 96,5 minutos.

O engajamento por câmera e microfone foi nulo (0,0\%), indicando ausência de interação audiovisual. Houve 13 trocas de aba/janela e intensidade de multitarefa de 21,0\%, resultando em fragmentação de foco de 240,0\%, o que reforça a alternância frequente entre a aula e outras atividades digitais.

A distribuição temporal do engajamento, mostrada na Figura~\ref{fig:aml_quartil}, apresenta padrão irregular, com pico às 09:50 e queda rápida para 50\% do engajamento máximo nos primeiros 20 minutos. Isso sugere perda precoce de motivação e possível recuperação de interesse apenas no final da sessão.

\begin{figure}[H]
    \centering
    \begin{minipage}{0.9\textwidth}
        \centering
        \includegraphics[width=.9\textwidth]{espeon/analysis.aml.quartil.png}
        \caption{Distribuição Temporal - Análise e Modelagem de Sistemas}
        \label{fig:aml_quartil}
    \end{minipage}\hfill
\end{figure}

Os \textit{scores} gerais reforçam esse quadro: \textit{score} de engajamento de 44,4\%, saúde da atenção de 91,0\% e risco de distração de 26,5\%. Apesar de momentos pontuais de foco qualificado, o padrão predominante foi de atenção instável e baixa interação.
